\documentclass[12pt]{article}
\usepackage[utf8]{inputenc}
\usepackage[T1]{fontenc}
\usepackage{hyperref}
\title{Auslan Language Resources for K-12 Curriculum Development}
\author{Compiled by Claude Sonnet 5.5 (claude-sonnet-5-5) for Andr\'as Kornai}
\date{}
\begin{document}
\maketitle
\section{Introduction}
Auslan (ISO 639-3 asf) is the sign language of the Australian Deaf community. Curriculum material found is state-based; no national hub was found.
\section{Primary Educational Resources}
\begin{itemize}
\item Victorian Curriculum F--10 Auslan resources (VCAA): \url{https://www.vcaa.vic.edu.au/curriculum/foundation-10/resources/languages/auslan/Pages/default.aspx}
\item NSW Auslan K--10 Syllabus support: \url{https://education.nsw.gov.au/teaching-and-learning/curriculum/languages/planning-programming-and-assessing-languages-7-10/auslan}
\item Languages NSW note on offering Auslan in 2026: \url{https://www.languagesnsw.com/post/is-your-school-offering-auslan-in-2026}
\item Twinkl commercial Auslan pack: \url{https://www.twinkl.com/resource/auslan-places-at-school-pack-au-al-1686539494}
\end{itemize}
\section{Linguistic and Reference Materials}
\begin{itemize}
\item Auslan Signbank dictionary: \url{https://auslan.org.au/}
\item Signbank and the Auslan Corpus (Monash): \url{https://www.monash.edu/arts/languages-literatures-cultures-linguistics/signbank-and-the-auslan-corpus}
\end{itemize}
\section{Religious and Community Resources}
None were found in this search.
\section{Digital and Audio Resources}
Signbank offers signing videos, searchable by topic including education.
\section{Curriculum Development Context}
Auslan is part of the Australian curriculum languages offer; the Victorian resources are aligned to Version 1.0 and may be superseded by Version 2.0, and NSW implementation begins in 2026.
\end{document}